\chapterHeading{Number theory}

\sectionHeading{Sieve of Eratosthenes}
\begin{lstlisting}[language=C++]
const int N = 1e7 + 5;
bitset<N> composite;
vector<int> primes;

void sieve() {
    composite[0] = composite[1] = 1;
    for (int i = 2; i < N; i++) {
        if (composite[i]) continue;
        primes.push_back(i);
        for (long long j = 1LL * i * i; j < N; j += i) composite[j] = 1;
    }
}
\end{lstlisting}

\sectionHeading{Segmented sieve (primes in [a, b])}
\begin{lstlisting}[language=C++]
vector<long long> segmentedSieve(long long a, long long b) {
    // assumes sieve() has already filled primes up to sqrt(b)
    vector<bool> isPrime(b - a + 1, true);
    if (a == 1) isPrime[0] = false;
    for (int p : primes) {
        if (1LL * p * p > b) break;
        long long start = max(1LL * p * p, (a + p - 1) / p * p);
        for (long long j = start; j <= b; j += p) isPrime[j - a] = false;
    }
    vector<long long> res;
    for (long long i = a; i <= b; i++)
        if (isPrime[i - a]) res.push_back(i);
    return res;
}
\end{lstlisting}

\sectionHeading{Number of divisors}
\begin{lstlisting}[language=C++]
const int N = 1e6 + 5;
int divisors[N];

void precomputeDivisors() {
    for (int i = 1; i < N; i++)
        for (int j = i; j < N; j += i) divisors[j]++;
}
\end{lstlisting}

\sectionHeading{Sum of divisors}
\begin{lstlisting}[language=C++]
const int N = 1e6 + 5;
long long divSum[N];

void precomputeDivisorSum() {
    for (int i = 1; i < N; i++)
        for (int j = i; j < N; j += i) divSum[j] += i;
}
\end{lstlisting}

\sectionHeading{Euler's totient (phi)}
\begin{lstlisting}[language=C++]
const int N = 1e6 + 5;
int phi[N];

void precomputePhi() {                 // O(N log log N)
    for (int i = 1; i < N; i++) phi[i] = i;
    for (int i = 2; i < N; i++)
        if (phi[i] == i)
            for (int j = i; j < N; j += i) phi[j] -= phi[j] / i;
}

int eulerPhi(int n) {                  // single value, O(sqrt n)
    int res = n;
    for (long long i = 2; i * i <= n; i++)
        if (n % i == 0) {
            while (n % i == 0) n /= i;
            res -= res / i;
        }
    if (n > 1) res -= res / n;
    return res;
}
\end{lstlisting}

\sectionHeading{Big mod (modular exponentiation)}
\begin{lstlisting}[language=C++]
long long bigMod(long long a, long long b, long long mod) {
    long long ans = 1;
    a %= mod;
    while (b) {
        if (b & 1) ans = ans * a % mod;
        a = a * a % mod;
        b >>= 1;
    }
    return ans;
}

// a, b, mod up to 1e18: use 128-bit multiplication
long long mulMod(long long a, long long b, long long m) {
    return (__int128)a * b % m;
}
long long bigMod128(long long a, long long b, long long m) {
    long long ans = 1;
    a %= m;
    while (b) {
        if (b & 1) ans = mulMod(ans, a, m);
        a = mulMod(a, a, m);
        b >>= 1;
    }
    return ans;
}
\end{lstlisting}

\sectionHeading{nCr}
\begin{lstlisting}[language=C++]
const int MX = 1e6 + 5;
long long fact[MX], invFact[MX];

void precomputeFact(long long mod) {
    fact[0] = 1;
    for (int i = 1; i < MX; i++) fact[i] = fact[i - 1] * i % mod;
    invFact[MX - 1] = bigMod(fact[MX - 1], mod - 2, mod);
    for (int i = MX - 1; i > 0; i--) invFact[i - 1] = invFact[i] * i % mod;
}

long long nCr(int n, int r, long long mod) {       // mod prime
    if (r < 0 || r > n) return 0;
    return fact[n] * invFact[r] % mod * invFact[n - r] % mod;
}

long long nCrNoMod(int n, int r) {                 // O(r), no mod
    if (r < 0 || r > n) return 0;
    r = min(r, n - r);
    long long res = 1;
    for (int i = 1; i <= r; i++) res = res * (n - r + i) / i;
    return res;
}
\end{lstlisting}

\sectionHeading{Josephus problem}
\begin{lstlisting}[language=C++]
int josephus(int n, int k) {           // 1-based survivor, O(n)
    int res = 0;
    for (int i = 1; i <= n; i++) res = (res + k) % i;
    return res + 1;
}
\end{lstlisting}
